\begin{pastexam}[label={px:ntust-113-ice-03}]{臺科大資工所 113 學年度 第 I-3 題}
（複選）

I.\quad \underline{(35\%) Multiple choice questions.}

(5\%) Select all the correct statements for finding the shortest paths from a given node in a graph.
\begin{enumerate}[label=(\alph*)]
\item Dijkstra's algorithm can be applied if the graph is undirected without negative weights.
\item Dijkstra's algorithm can be applied if the graph is directed with negative weights but without
negative cycles.
\item Dijkstra's algorithm can be applied if the graph is undirected with negative cycles.
\item Bellman-ford algorithm can be applied if the graph is undirected without negative weights.
\item Bellman-ford algorithm can indicate that no solution exists if the graph is directed with
negative cycles.
\item The answers must be the same for every graph that can successfully find shortest path
solutions using both Dijkstra's and Bellman-ford algorithms.
\end{enumerate}
\end{pastexam}

\begin{pastsoln}
本書的答案是 (a)、(d)、(e)、(f)。

無向圖的每條邊 $\{u, v\}$ 都可以看成兩條有向邊 $u \to v$、$v \to u$，權重相同；有向圖的單源最短
路徑演算法都可以這樣套用到無向圖上。
\begin{itemize}
\item \textbf{(a) 對。} 換成有向圖後權重仍然都 $\ge 0$，滿足 Dijkstra 的條件（資結書定理~25.6）。
\item \textbf{(b) 錯。} Dijkstra 要求權重都 $\ge 0$，沒有負環但有負邊時就可能算錯（資結書
第~25.2.5~節）。 反例：$s \to a$ 權重 $2$、$s \to b$ 權重 $3$、$b \to a$ 權重 $-2$。 Dijkstra 先確定
$d[a] = 2$，之後才取出 $b$，但這時 $a$ 已經確定，不會再改；正確答案 $\delta(s, a) = 1$。
\item \textbf{(c) 錯。} 有負環時，經過負環可以無限變小，最短路徑根本不存在，任何演算法都給不出
答案；Dijkstra 連負邊都處理不了。 （無向圖只要有一條負邊 $\{u, v\}$，$u \to v \to u$ 就是負環。）
\item \textbf{(d) 對。} Bellman--Ford 只要求從起點走得到的地方沒有負環（資結書定理~25.7）。 無向圖
沒有負權重，換成兩條有向邊後也沒有負環，可以用。
\item \textbf{(e) 對。} Bellman--Ford 做完 $|V| - 1$ 輪放鬆後再檢查一輪，還能放鬆就回傳 \textsc{false}，
表示有從起點走得到的負環、最短路徑不存在（資結書第~25.3.4~節）。 若負環從起點走不到，那些點的
最短路徑仍然存在，它也正確回傳；題目說「can indicate」，成立。
\item \textbf{(f) 對。} 兩者都成功時，算出的是同一個圖的最短路徑距離 $\delta(s, v)$，而最短距離是
唯一的，所以答案一樣。 （最短路徑本身可能有好幾條一樣短的，兩個演算法回報的路徑可以不同，但
距離一定相同；題目的「answers」指最短路徑的解，也就是距離。）
\end{itemize}
\end{pastsoln}
